% Propietario conservado: /Users/ruben/Documents/ChatGPT/jueces y controles/REVISION_ARGUMENTAL_APERTURAS_CIERRES_20260915/delivery/ARTICLE_V_ES_ARGUMENTAL_20260915/manuscrito/sections/40_exclusion_pauli.tex
% Recuperación documental para el artículo V; RESULTADO_RECUPERADO.
% Propietario: /Users/ruben/Documents/New project/output/REVISION_EDITORIAL_HMT_MD_20260905/01_FUENTES/INTEGRAL/tree/New project/output/ENSAMBLADO_CAUSAL_RESTAURADO_HMT_MD_20260905/fuente/manuscrito/integracion_83/parte_iv/body/B015_ch56_pauli_completacion_fermionica_7549b83967f5_04d_ocupacion_estadistica.tex
% Líneas de origen: 1--48.
% REV02: aplicación estadística; prueba común de idempotencia reunida en 20.
% Recupera además los recuentos y ocupaciones antes anticipados en 20.
% Correspondencia editorial: gestion/CONSOLIDACION_PAULI_FOCK_REV02.json.
\section{Application statistique des spectres d’occupation}
\label{v:v:rec:chap:ocupacion-estadistica}
\index{Pauli, principe d’exclusion de}\index{CAR}
\index{Fermi--Dirac}\index{Bose--Einstein}

Les réalisations nulle et matérielle fournissent le support géométrique de
l’occupation. La complétion graduée de ce support détermine ses espaces
multiparticulaires et ses opérateurs ; le
théorème~\ref{v:v:rec:fock_gibbs:statement:01} établit leur transport par
raffinement. Cette étape utilise ces résultats pour caractériser les
configurations admissibles, dénombrer les états et préparer l’évaluation thermique.
La construction algébrique, le dénombrement fini et l’état d’équilibre
conservent ainsi leurs fonctions et leurs prémisses respectives.

\subsection{Restriction spectrale et configurations admissibles}
\label{v:v:rec:pauli:section:02}

Pour un mode \(i\), nous écrivons \(n_i=N_i=a_i^\dagger a_i\). Le
théorème~\ref{v:v:rec:thm:principio-exclusion-pauli}, dont la preuve complète figure dans
\eqref{v:v:pauli:eq:prueba-idempotencia}, établit
\(n_i^2=n_i\) et \(\Spec(n_i)\subseteq\{0,1\}\). Le calcul utilise
uniquement CAR et la relation d’adjonction ; il s’applique donc également
à toute représentation de ces relations. La nilpotence
\((a_i^\dagger)^2=0\) empêche la double occupation fermionique du même mode.
Par conséquent, une configuration fermionique de \(g\) modes est une
famille \((n_1,\ldots,n_g)\in\{0,1\}^g\) ; dans le secteur bosonique, les
occupations appartiennent à \(\mathbb N_0^g\).

La représentation fermionique utilise l’algèbre CAR construite à partir de la
complétion extérieure. L’identification relativiste de
spin--statistique conserve, en outre, la compatibilité de la représentation
des rotations avec la parité et les conditions de localité exposées dans
les sections~\ref{v:v:pauli:car-construccion} et
\ref{v:v:rec:completaciones:section:05}. La périodicité \(4\pi\) et
l’exclusion sont coordonnées par le même caractère de feuillet, avec ces
conditions de représentation explicites.

\subsection{Dimensions des secteurs d’occupation fixe}
\label{v:v:rec:completaciones:section:07}

S’il y a \(g\) modes et \(N\) particules, les dimensions des secteurs sont
\begin{equation}
 \Omega_F(g,N)=\binom gN,
 \qquad
 \Omega_B(g,N)=\binom{g+N-1}{N}.
 \label{v:v:ocupacion:eq:dimensiones-sectores}
\end{equation}
La première formule dénombre des sous-ensembles de modes, avec une dimension nulle si
\(N>g\) ; la seconde, des compositions faibles de \(N\) en \(g\) parties.
Ces dénombrements sont des conséquences exactes des deux complétions et
s’obtiennent avant l’introduction de l’état thermique. La
proposition~\ref{v:v:rec:prop:cardinalidades-microcanonicas} applique les deux
dénombrements à l’écran spectral fini, en conservant ses niveaux,
ses multiplicités et ses contraintes d’énergie et de population.

\subsection{Lecture thermique des complétions graduées}
\label{v:v:ocupacion:lectura-termica}

Avec un Hamiltonien diagonal \(H\), une énergie modale \(\epsilon\), un potentiel
chimique \(\mu\) et un paramètre thermique \(\beta>0\), les occupations moyennes
des deux secteurs sont
\begin{equation}
 \bar n_{\rm par}(\epsilon)
 =\frac{1}{e^{\beta(\epsilon-\mu)}-1},
 \qquad
 \bar n_{\rm impar}(\epsilon)
 =\frac{1}{e^{\beta(\epsilon-\mu)}+1}.
 \label{v:v:ocupacion:eq:leyes-por-modo}
\end{equation}
Dans le secteur bosonique, on exige \(\mu<\inf\operatorname{Spec}H\).
Ces expressions sont la spécialisation par mode du
théorème~\ref{v:v:rec:thm:factorizacion-gran-canonica-rev3}, dont la preuve
construit la trace grand canonique à partir des occupations admissibles.
Les formules de Bose--Einstein et de Fermi--Dirac sont donc les
évaluations thermodynamiques des complétions symétrique et extérieure
sélectionnées par le transport. La même section obtient leurs expressions
au moyen de la variation de l’entropie sous les contraintes d’énergie et de
nombre ; le théorème~\ref{v:v:rec:thm:limite-clasico-fd-be-rev2} démontre la
limite classique commune avec une borne explicite de l’erreur.
